% Options for packages loaded elsewhere
% Options for packages loaded elsewhere
\PassOptionsToPackage{unicode}{hyperref}
\PassOptionsToPackage{hyphens}{url}
%
\documentclass[
  english,
  russian,
  12pt,
  a4paper,
  DIV=11,
  numbers=noendperiod]{scrreprt}
\usepackage{xcolor}
\usepackage{amsmath,amssymb}
\setcounter{secnumdepth}{5}
\usepackage{iftex}
\ifPDFTeX
  \usepackage[T1]{fontenc}
  \usepackage[utf8]{inputenc}
  \usepackage{textcomp} % provide euro and other symbols
\else % if luatex or xetex
  \usepackage{unicode-math} % this also loads fontspec
  \defaultfontfeatures{Scale=MatchLowercase}
  \defaultfontfeatures[\rmfamily]{Ligatures=TeX,Scale=1}
\fi
\usepackage{lmodern}
\ifPDFTeX\else
  % xetex/luatex font selection
\fi
% Use upquote if available, for straight quotes in verbatim environments
\IfFileExists{upquote.sty}{\usepackage{upquote}}{}
\IfFileExists{microtype.sty}{% use microtype if available
  \usepackage[]{microtype}
  \UseMicrotypeSet[protrusion]{basicmath} % disable protrusion for tt fonts
}{}
\usepackage{setspace}
% Make \paragraph and \subparagraph free-standing
\makeatletter
\ifx\paragraph\undefined\else
  \let\oldparagraph\paragraph
  \renewcommand{\paragraph}{
    \@ifstar
      \xxxParagraphStar
      \xxxParagraphNoStar
  }
  \newcommand{\xxxParagraphStar}[1]{\oldparagraph*{#1}\mbox{}}
  \newcommand{\xxxParagraphNoStar}[1]{\oldparagraph{#1}\mbox{}}
\fi
\ifx\subparagraph\undefined\else
  \let\oldsubparagraph\subparagraph
  \renewcommand{\subparagraph}{
    \@ifstar
      \xxxSubParagraphStar
      \xxxSubParagraphNoStar
  }
  \newcommand{\xxxSubParagraphStar}[1]{\oldsubparagraph*{#1}\mbox{}}
  \newcommand{\xxxSubParagraphNoStar}[1]{\oldsubparagraph{#1}\mbox{}}
\fi
\makeatother


\usepackage{longtable,booktabs,array}
\usepackage{calc} % for calculating minipage widths
% Correct order of tables after \paragraph or \subparagraph
\usepackage{etoolbox}
\makeatletter
\patchcmd\longtable{\par}{\if@noskipsec\mbox{}\fi\par}{}{}
\makeatother
% Allow footnotes in longtable head/foot
\IfFileExists{footnotehyper.sty}{\usepackage{footnotehyper}}{\usepackage{footnote}}
\makesavenoteenv{longtable}
\usepackage{graphicx}
\makeatletter
\newsavebox\pandoc@box
\newcommand*\pandocbounded[1]{% scales image to fit in text height/width
  \sbox\pandoc@box{#1}%
  \Gscale@div\@tempa{\textheight}{\dimexpr\ht\pandoc@box+\dp\pandoc@box\relax}%
  \Gscale@div\@tempb{\linewidth}{\wd\pandoc@box}%
  \ifdim\@tempb\p@<\@tempa\p@\let\@tempa\@tempb\fi% select the smaller of both
  \ifdim\@tempa\p@<\p@\scalebox{\@tempa}{\usebox\pandoc@box}%
  \else\usebox{\pandoc@box}%
  \fi%
}
% Set default figure placement to htbp
\def\fps@figure{htbp}
\makeatother



\ifLuaTeX
\usepackage[bidi=basic,provide=*]{babel}
\else
\usepackage[bidi=default,provide=*]{babel}
\fi
% get rid of language-specific shorthands (see #6817):
\let\LanguageShortHands\languageshorthands
\def\languageshorthands#1{}


\setlength{\emergencystretch}{3em} % prevent overfull lines

\providecommand{\tightlist}{%
  \setlength{\itemsep}{0pt}\setlength{\parskip}{0pt}}



 
\usepackage[style=gost-numeric,backend=biber,langhook=extras,autolang=other*]{biblatex}
\addbibresource{bib/cite.bib}

\usepackage[]{csquotes}

\usepackage{indentfirst}
% \usepackage{float}
% \floatplacement{figure}{H}
\IfFileExists{plex-otf.sty}{
  %% Full TeXlive
  \usepackage[math,RM={Scale=0.94},SS={Scale=0.94},SScon={Scale=0.94},TT={Scale=MatchLowercase,FakeStretch=0.9},DefaultFeatures={Ligatures=Common}]{plex-otf}
}{
  %% TinyTeX
  \usepackage{libertine}
}

%%% Расширенная математика

\usepackage{amsmath}
\usepackage{amssymb}
\usepackage{amscd}

\usepackage{mathtools}
\mathtoolsset{
  showonlyrefs=true,
  mathic=true,
}

\allowdisplaybreaks

%%%% Calligraphic fonts for physicists
\usepackage{mathrsfs}

%%% Local Variables:
%%% mode: latex
%%% coding: utf-8-unix
%%% End:
\RequirePackage{scrhack}
\RequirePackage{geometry}

% \geometry{includehead}
\geometry{includefoot}
\geometry{left=30mm,right=15mm,top=20mm,bottom=20mm}
\geometry{bindingoffset=0mm}
\geometry{marginparwidth=0dd,marginparsep=0dd}
\geometry{heightrounded}

%%% Titles

\providecommand{\DivFontShape}{\bfseries}
\providecommand{\DivFontSize}{}
\providecommand{\DivPosition}{}
%%
\providecommand{\partFontShape}{\DivFontShape}
\providecommand{\partFontSize}{\huge}
\providecommand{\partPosition}{}
%%
\providecommand{\chapterFontShape}{\DivFontShape}
\providecommand{\chapterFontSize}{\Large}
\providecommand{\chapterPosition}{}
%%
\providecommand{\sectionFontShape}{\DivFontShape}
\providecommand{\sectionFontSize}{\large}
\providecommand{\sectionPosition}{\center}
%%
\providecommand{\subsectionFontShape}{\DivFontShape}
\providecommand{\subsectionFontSize}{\large}
\providecommand{\subsectionPosition}{\center}
%%
\providecommand{\subsubsectionFontShape}{\DivFontShape}
\providecommand{\subsubsectionFontSize}{\normalsize}
\providecommand{\subsubsectionPosition}{\center}
%%
\providecommand{\paragraphFontShape}{\DivFontShape}
\providecommand{\paragraphFontSize}{\normalsize}
\providecommand{\paragraphPosition}{}
%%
\providecommand{\subparagraphFontShape}{\DivFontShape}
\providecommand{\subparagraphFontSize}{\normalsize}
\providecommand{\subparagraphPosition}{}
%%
\providecommand{\contentsnameFontShape}{\DivFontShape}
\providecommand{\contentsnameFontSize}{\Large}
\providecommand{\contentsnamePosition}{\flushleft}

\KOMAoptions{toc=chapterentrydotfill}
\addtokomafont{chapterentrypagenumber}{\normalfont\chapterFontShape}

% \RequirePackage{tocloft}
\RequirePackage[titletoc]{appendix}

% \def\@pnumwidth{1.55em}
% \def\@tocrmarg {2.55em}
% \def\@dotsep{4.5}
\setcounter{tocdepth}{2}

%%% Captions

\RequirePackage{caption}

\captionsetup{font={small}}
\captionsetup{textfont={bf},labelfont={bf}}

% \DeclareCaptionFormat{tableRight}{\hfill#1\newline#2#3\par}
\DeclareCaptionFormat{gost1.5}{\hfill#1\newline#2#3\par}

\captionsetup[figure]{labelsep=period,position=bottom,justification=centering}

% \captionsetup[table]{format=tableRight,labelsep=none,position=top}
\captionsetup[table]{format=gost1.5,labelsep=none,position=top,justification=centering}

%\AtBeginDocument{%
% \RequirePackage{floatrow}
%\makeatletter
%\@ifpackageloaded{floatrow}{%
%    \floatsetup[table]{style=Plaintop}
%    }{}
%\makeatother
%}

%%%  Chapter

\KOMAoptions{headings=standardclasses,
  headings=big,
chapterprefix=false}
\renewcommand*\raggedchapter{\centering}
\renewcommand*\raggedsection{\centering}
\RedeclareSectionCommand[beforeskip=0pt,afterskip=1\baselineskip]{chapter}
% \setkomafont{chapterprefix}{\normalsize\mdseries}

%%% Russian

\usepackage{indentfirst}

%%% Lists

\usepackage{enumitem}
\setlist{nosep}
\setlist{noitemsep}
\setlist[enumerate]{labelsep=*, leftmargin=*}
\setlist[itemize]{leftmargin=*,topsep=0pt}
\setlist[itemize,1]{label=---~}
\setlist[itemize,2]{label=---~}
\setlist[itemize,3]{label=---~}

%%% Local Variables:
%%% mode: latex
%%% coding: utf-8-unix
%%% End:
\KOMAoption{captions}{tableheading}
\makeatletter
\@ifpackageloaded{bookmark}{}{\usepackage{bookmark}}
\makeatother
\makeatletter
\@ifpackageloaded{caption}{}{\usepackage{caption}}
\AtBeginDocument{%
\ifdefined\contentsname
  \renewcommand*\contentsname{Содержание}
\else
  \newcommand\contentsname{Содержание}
\fi
\ifdefined\listfigurename
  \renewcommand*\listfigurename{Список иллюстраций}
\else
  \newcommand\listfigurename{Список иллюстраций}
\fi
\ifdefined\listtablename
  \renewcommand*\listtablename{Список таблиц}
\else
  \newcommand\listtablename{Список таблиц}
\fi
\ifdefined\figurename
  \renewcommand*\figurename{Рисунок}
\else
  \newcommand\figurename{Рисунок}
\fi
\ifdefined\tablename
  \renewcommand*\tablename{Таблица}
\else
  \newcommand\tablename{Таблица}
\fi
}
\@ifpackageloaded{float}{}{\usepackage{float}}
\floatstyle{ruled}
\@ifundefined{c@chapter}{\newfloat{codelisting}{h}{lop}}{\newfloat{codelisting}{h}{lop}[chapter]}
\floatname{codelisting}{Список}
\newcommand*\listoflistings{\listof{codelisting}{Листинги}}
\makeatother
\makeatletter
\makeatother
\makeatletter
\@ifpackageloaded{caption}{}{\usepackage{caption}}
\@ifpackageloaded{subcaption}{}{\usepackage{subcaption}}
\makeatother
	    \csundef{listoflistings}
	    
\makeatletter
\@ifpackageloaded{minted}{}{\usepackage[chapter, ]{minted}}
\makeatother
	    \usemintedstyle{emacs}
	    \setminted{mathescape=true}
	    \setminted{autogobble=true}
	    \setminted{breaklines=true}
	    \setminted{breakanywhere=true}
	    \setminted{frame=lines}
	    % \setminted{bgcolor=lightgray}
	    %\setminted{linenos=true,
	    %  numberblanklines=false,
	    %}  
	    
\usepackage{bookmark}
\IfFileExists{xurl.sty}{\usepackage{xurl}}{} % add URL line breaks if available
\urlstyle{same}
\hypersetup{
  pdfauthor={Кулябов Д. С.},
  pdflang={ru},
  hidelinks,
  pdfcreator={LaTeX via pandoc}}


\author{Кулябов Д. С.}
\date{}
\begin{document}

\renewcommand*\contentsname{Содержание}
{
\setcounter{tocdepth}{1}
\tableofcontents
}

\setstretch{1.5}
\bookmarksetup{startatroot}

\chapter*{}\label{section}

\markboth{}{}

\bookmarksetup{startatroot}

\chapter*{Введение}\label{ux432ux432ux435ux434ux435ux43dux438ux435}
\addcontentsline{toc}{chapter}{Введение}

\markboth{Введение}{Введение}

Данная работа посвящена исследованию средств натурного моделирования на
примере распространения компьютерного вируса. Эта работа направлена на
анализ и сравнение трех эмуляторов: Kathara, Mininnet и Containerlab.
Для изучения характеристик средств моделирования рассматриваются
эпидемиалогическая модель SEIRS-NIMFA и компьютерный червь Морриса.

\section*{Актуальность
работы}\label{ux430ux43aux442ux443ux430ux43bux44cux43dux43eux441ux442ux44c-ux440ux430ux431ux43eux442ux44b}
\addcontentsline{toc}{section}{Актуальность работы}

\markright{Актуальность работы}

Актуальность данной работы обусловлена отсутствием единого исследования,
рассматривающего все три средства моделирования. Эта научная работа
напраавлена на заполнение данного пробела знаний эмпирическим путем.

С помощью проведения эксперимента по распространению компьютерного
вируса будут выявлены отличительные характеристики эмуляторов, что
поможет сформулировать вывод о релевантности их использования для
решения исследовательских задач и даст понимание о возможностях и
ограничениях рассматриваемых средств натурного моделирования.

\section*{Цель
работы}\label{ux446ux435ux43bux44c-ux440ux430ux431ux43eux442ux44b}
\addcontentsline{toc}{section}{Цель работы}

\markright{Цель работы}

Целью выпускной квалификационной работы является исследование средств
натурного моделирования путем распространения компьютерного вируса.

\section*{Задачи
работы}\label{ux437ux430ux434ux430ux447ux438-ux440ux430ux431ux43eux442ux44b}
\addcontentsline{toc}{section}{Задачи работы}

\markright{Задачи работы}

Основными задачами работы являются:

\begin{enumerate}
\def\labelenumi{\arabic{enumi}.}
\tightlist
\item
  Аназил средств натурного моделирования: Kathara, Mininet и
  Containerlab.
\item
  Теоретическое изучение эпидемиалогической модели SEIRS-NIMFA и
  компьютерного червя Морриса.
\item
  Рассмотрение алгоритмов для моделирования распространения
  компьютерного вируса в выбранных эмуляторах.
\item
  Программная реализация моделируемых экспериментов в средствах
  моделирования.
\end{enumerate}

\section*{Методы
исследования}\label{ux43cux435ux442ux43eux434ux44b-ux438ux441ux441ux43bux435ux434ux43eux432ux430ux43dux438ux44f}
\addcontentsline{toc}{section}{Методы исследования}

\markright{Методы исследования}

Методами исследования является изучение и анализ научной литературы по
теме, а также создание и проведение воспроизводимого эксперимента
распространения вируса.

\section*{Структура
работы}\label{ux441ux442ux440ux443ux43aux442ux443ux440ux430-ux440ux430ux431ux43eux442ux44b}
\addcontentsline{toc}{section}{Структура работы}

\markright{Структура работы}

Работа состоит из введения, трех разделов, заключения и списка
используемой литературы.

Во введении раскрыта суть рассматриваемой проблемы, обоснована
актуальность темы, сформулированы цели и задачи, определены методы
исследования.

В первом разделе \ldots{}

Во втором разделе теоретически рассматриваются эпидемиалогические модели
и червь Морриса.

В третьем разделе описаны параметры установки средств моделирования и
выполнена программная реализация эксперимента. Проведен сравнительный
анализ характеристик эмуляторов и получены результаты для обсуждения.

В заключении подведены общие итоги выпускной квалификационной работы,
изложены основные выводы.

\bookmarksetup{startatroot}

\chapter{Теоретическая часть курсовой
работы}\label{ux442ux435ux43eux440ux435ux442ux438ux447ux435ux441ux43aux430ux44f-ux447ux430ux441ux442ux44c-ux43aux443ux440ux441ux43eux432ux43eux439-ux440ux430ux431ux43eux442ux44b}

\section{Компартментные
модели}\label{ux43aux43eux43cux43fux430ux440ux442ux43cux435ux43dux442ux43dux44bux435-ux43cux43eux434ux435ux43bux438}

Компартментные модели - математическая модель, которая описывает
процессы взаимодействия различных субъектов (например, особей) из разных
групп (компартментов) в течение времени. Каждый субъект принадлежит
одной группе, при этом внутри каждой группы субъекты неразличимы.

Одной из простейших компартментных моделей является модель
распространения эпидемий SIR. Эта модель была создана Уильямом Кермаком
и Андерсоном МакКендриком в 1927 году. Все особи популяции в данном
модели делсятся на три группы:

\begin{itemize}
\tightlist
\item
  S (susceptible) - здоровые особи, подверженные заболеванию;
\item
  I (infectious) - зараженные особи, распространяющие болезнь;
\item
  R (recovered) - переболевшие особи, приобретшие иммунитет.
\end{itemize}

Однако, во многих болезнях есть икубационных период, во время которого
особь заражена, но не пока заразна. Так появилась модель SEIR, где
присутвует четвертая группа E (exposed) - особи, находящиеся в латентном
периоде заболевания. Такая модель повышает точность моделирования
инфекционных заболеваний, например, COVID-19.

Мы будем рассматривать замкнутую популяцию, где отсутствуют процессы
рождаемости и смертности. Тогда можно описать эволюцию особи следующей
диаграммой:

\[
S \xrightarrow{\beta} E \xrightarrow{\sigma} I \xrightarrow{\gamma} R
\]

Модель описывается системой дифференциальных уравнений:

\[
\begin{aligned}
\frac{dS}{dt} &= -\frac{\beta SI}{N}, \\
\frac{dE}{dt} &= \frac{\beta SI}{N} - \sigma E, \\
\frac{dI}{dt} &= \sigma E - \gamma I, \\
\frac{dR}{dt} &= \gamma I, \\
N &= S(t) + E(t) + I(t) + R(t),
\end{aligned}
\]

где - \(\beta\) - коэффициент заражения (вероятность того, что контакт
между восприимчивым и заражённым приводит к новому заражению)
(\(S \rightarrow E\)); - \(\sigma\) - коэффициент инкубационного
перехода (вероятность того, что зараженный индивид становится заразным)
(\(E \rightarrow I\)); - \(\gamma\) - коэффициент выздоровления
(вероятность того, что заражённый индивид выздоравливает)
(\(I \rightarrow R\)).

Еще одной модификацией модели SIR является модель SEIRS. В жизни у
определенной части переболевших особей иммунитет со временем ослабевает.
Модель SEIRS допускает перехо особей из состояния выздоровевших в
восприимчивое состояние.

Диаграмма эволюции особи выглядит следующим образом (рис.
\ref{fig-seirs_diagram}).

\begin{figure}

\centering{

\includegraphics[width=0.7\linewidth,height=\textheight,keepaspectratio]{text/../image/seirs_diagram.png}

}

\caption{\label{fig-seirs_diagram}Диаграмма эволюции особи в модели
SEIRS}

\end{figure}%

Система дифференциальных уравнений для модели SEIRS имеет вид:

\[
\begin{aligned}
\frac{dS}{dt} &= -\frac{\beta S I}{N} + \xi R, \\
\frac{dE}{dt} &= \frac{\beta S I}{N} - \sigma E, \\
\frac{dI}{dt} &= \sigma E - \gamma I, \\
\frac{dR}{dt} &= \gamma I - \xi R \\
N &= S(t) + E(t) + I(t) + R(t),
\end{aligned}
\]

где - \(\xi\) - вероятность потери иммунитета (вероятность того, что
выздоровевший индивид со временем возвращается в категорию
восприимчивых) (\(R \rightarrow S\)).

Если приток восприимчивых в популяцию достаточно велик, система в
установившемся состоянии переходит в эндемическое равновесие (устойчивое
состояние системы, в котором инфекция сохраняется на постоянном уровне:
не исчезает и не угасает), сопровождаемое затухающими колебаниями
численности заболевших.

Классические модели SIR/SEIR описывают усреднённую динамику
инфекционного процесса в популяции и не учитывают структуру
взаимодействий между её особями. В реальных же системах --- например, в
сетях Интернета вещей (IoT) --- заражение распространяется по связям
между конкретными устройствами, аналогично тому, как в популяции особи
контактируют не со всеми, а только с соседями.

Для учёта такого взаимодействия используется метод NIMFA (N-Intertwined
Mean-Field Approximation). В нём каждая особь из классической модели
рассматривается как отдельный узел графа, а вероятность её заражения
определяется степенью заражённости соседних узлов.

В данной работе рассматривается модель SEIRS-NIMFA, которая объединяет
принципы SEIRS-модели и сетевого приближения NIMFA. В отличие от
усреднённых моделей, здесь анализируется общая эволюция сети (количество
элементов в каждом компартменте) и индивидуальное состояние каждого
узла.

Вместо того чтобы рассматривать фактическое состояние каждого устройства
в сети и все возможные комбинации состояний для N узлов на каждом
временном шаге, модель учитывает вероятности для каждого компартмента.

Метод NIMFA описывает вероятности заражения \(p_i(t)\) для каждого узла
\(i\) во времени с помощью дифференциального уравнения:

\[
\frac{dp_i}{dt} = -\delta p_i(t) + (1 - p_i(t))\beta \sum_{j=1}^{N} a_{ij} p_j(t).
\]

где: - \(p_i(t)\) --- вероятность того, что узел i заражён в момент
времени t; - \(\beta\) --- коэффициент заражения; - \(\delta\) ---
коэффициент восстановления (скорость устранения заражения); - \(a_{ij}\)
--- элемент матрицы смежности, определяющий связь между узлами i и j
(т.е., если вершины графа связаны \(a_{ij} = 1\), иначе - 0).

Таким образом, вероятность заражения конкретного узла зависит от
состояния его соседей и от параметров распространения вредоносного ПО.

Для каждого устройства \(i\) вводятся вероятности нахождения в четырёх
состояниях:

\[
x_i(t),\; w_i(t),\; y_i(t),\; z_i(t),
\]

соответствующие компартментам S, E, I, R, при этом выполняется
нормировка:

\[
x_i(t) + w_i(t) + y_i(t) + z_i(t) = 1.
\]

Эволюцию особи можно описать следующей диагрмаммой (рис.
\ref{fig-seirs_nimfa_diagram}).

\begin{figure}

\centering{

\includegraphics[width=0.7\linewidth,height=\textheight,keepaspectratio]{text/../image/seirs_nimfa_diagram.png}

}

\caption{\label{fig-seirs_nimfa_diagram}Диаграмма эволюции особи в
модели SEIRS-NIMFA}

\end{figure}%

Скорость заражения узла \(i\) зависит от вероятности заражённости его
соседей и выражается как:

\[
b_i(t) = \beta_i \sum_{j=1}^{N} a_{ij} y_j(t),
\] где \(a_{ij}\) --- элемент матрицы смежности, показывающий наличие
связи между узлами.

Модель SEIRS-NIMFA описывается следующей системой дифференциальных
уравнений:

\[
\begin{aligned}
\frac{dx_i}{dt} &= -b_i(t)x_i + \gamma_i z_i,\\
\frac{dw_i}{dt} &= b_i(t)x_i - \alpha_i w_i, \\
\frac{dy_i}{dt} &= \alpha_i w_i - \delta_i y_i,\\
\frac{dz_i}{dt} &= \delta_i y_i - \gamma_i z_i.
\end{aligned}
\]

где

\begin{itemize}
\tightlist
\item
  \(\beta_i\) --- вероятность заразиться при контакте с инфицированным
  соседом (\(S \rightarrow E\));
\item
  \(\alpha_i\) --- вероятность перехода из инкубационного периода в
  зараженное состояние (\(E \rightarrow I\));
\item
  \(\delta_i\) --- вероятность выздоровления (\(I \rightarrow R\));
\item
  \(\gamma_i\) --- вероятность потери иммунитета и возвращения в
  восприимчивое состояние (\(R \rightarrow S\)).
\end{itemize}

\bookmarksetup{startatroot}

\chapter{Теоретическая часть курсовой
работы}\label{ux442ux435ux43eux440ux435ux442ux438ux447ux435ux441ux43aux430ux44f-ux447ux430ux441ux442ux44c-ux43aux443ux440ux441ux43eux432ux43eux439-ux440ux430ux431ux43eux442ux44b-1}

Основная часть курсовой работы содержит две или три главы.

Название глав не должно повторять название темы, а названия разделов ---
названия глав.

В \emph{первой главе} курсовой работы определяются теоретические аспекты
темы. Здесь излагается теоретическая сущность решаемой задачи,
рассматриваются различные подходы к решению, дается их оценка,
обосновывается и излагается собственная точка зрения.

Сдесь же предполагается теоретическое описание исследуемой темы.

\begin{gather}
  \Dot{W}( t ) =
  \frac{1}{T\left( Q,t \right)}\vartheta
  \left(W_{\max} - W\left( t \right) \right) -
  \frac{W(t) W (t - T(Q,t))}{2T(t - T(Q,t))}
  p ( t - T( Q,t ) ),
  \label{eq:dotW}
  \\
  \Dot{Q}(t) =
  \begin{dcases}
    (1 - p(t)) \frac{W(t)}{T(Q,t)} N (t) - C,
    \quad
    Q(t) > 0
    \\
    \max [ (1 - p(t)) \frac{W (t)}{T(Q,t)} N(t) - C, 0 ],
    \quad
    Q(t) = 0
  \end{dcases}
  \label{eq:dotQ}
  \\
  \Dot{\Hat{Q}} (t) = - \omega_{q} C\Hat{Q} (t) + \omega_{q} C
  \Hat{Q}(t)
  \label{eq:dotQhat}
\end{gather}

Одним из разделов теоретической части является \emph{литературный
обзор}. В данном разделе описываются результаты предшественников.

Также, следует приводить табличные материалы. Например, в табл.
\ref{tbl-std-dir} приведено краткое описание стандартных каталогов Unix.

\begin{longtable}[]{@{}
  >{\raggedright\arraybackslash}p{(\linewidth - 2\tabcolsep) * \real{0.1014}}
  >{\raggedright\arraybackslash}p{(\linewidth - 2\tabcolsep) * \real{0.8986}}@{}}
\caption{Описание некоторых каталогов файловой системы GNU
Linux}\label{tbl-std-dir}\tabularnewline
\toprule\noalign{}
\begin{minipage}[b]{\linewidth}\raggedright
Имя каталога
\end{minipage} & \begin{minipage}[b]{\linewidth}\raggedright
Описание каталога
\end{minipage} \\
\midrule\noalign{}
\endfirsthead
\toprule\noalign{}
\begin{minipage}[b]{\linewidth}\raggedright
Имя каталога
\end{minipage} & \begin{minipage}[b]{\linewidth}\raggedright
Описание каталога
\end{minipage} \\
\midrule\noalign{}
\endhead
\bottomrule\noalign{}
\endlastfoot
/ & Корневая директория, содержащая всю файловую \\
/bin & Основные системные утилиты, необходимые как в
однопользовательском режиме, так и при обычной работе всем
пользователям \\
/etc & Общесистемные конфигурационные файлы и файлы конфигурации
установленных программ \\
/home & Содержит домашние директории пользователей, которые, в свою
очередь, содержат персональные настройки и данные пользователя \\
/media & Точки монтирования для сменных носителей \\
/root & Домашняя директория пользователя root \\
/tmp & Временные файлы \\
/usr & Вторичная иерархия для данных пользователя \\
\end{longtable}

Более подробно об используемых утилитах см. в
\autocite{gnu-doc_bash,newham_2005_bash,zarrelli_2017_bash,robbins_2013_bash,tannenbaum_arch-pc_ru,tannenbaum_modern-os_ru}.

\bookmarksetup{startatroot}

\chapter{Аналитическая часть курсовой
работы}\label{ux430ux43dux430ux43bux438ux442ux438ux447ux435ux441ux43aux430ux44f-ux447ux430ux441ux442ux44c-ux43aux443ux440ux441ux43eux432ux43eux439-ux440ux430ux431ux43eux442ux44b}

\section{Проектирование структуры и компонентов программного
комплекса}\label{ux43fux440ux43eux435ux43aux442ux438ux440ux43eux432ux430ux43dux438ux435-ux441ux442ux440ux443ux43aux442ux443ux440ux44b-ux438-ux43aux43eux43cux43fux43eux43dux435ux43dux442ux43eux432-ux43fux440ux43eux433ux440ux430ux43cux43cux43dux43eux433ux43e-ux43aux43eux43cux43fux43bux435ux43aux441ux430}

После уточнения интерфейса выполняется декомпозиция предметной области
задачи в соответствии с выбранной технологией, т. е. создается
структурная схема будущего продукта и описывается взаимодействие его
функциональных элементов.

Структурная схема --- схема, отражающая состав и взаимодействие по
управлению частей разрабатываемого продукта. При объектной декомпозиции
такими частями являются объекты, при структурной декомпозиции ---
подпрограммы.

Для тем, связанных с нечисловой обработкой данных, этот раздел также
должен содержать информационную модель системы, которая может быть
представлена функциональной схемой или диаграммой потоков данных.
Функциональная схема --- схема взаимодействия частей системы с описанием
информационных потоков, состава данных в потоках и указанием
используемых файлов и устройств.

Далее описывается проектирование компонентов в соответствии с выбранной
технологией. Для программы, использующей структурный подход к
программированию, в данном разделе приводятся обобщенные алгоритмы,
например, алгоритм основной программы и описывается межпрограммный
интерфейс подпрограмм.

Всего приводятся два--три из наиболее интересных алгоритмов, включая
алгоритм основной программы при структурном подходе.

Каждый алгоритм должен быть представлен:

\begin{itemize}
\tightlist
\item
  схемой алгоритма (алгоритм может быть представлен в псевдокоде или в
  виде блок-схем);
\item
  описанием процесса обработки данных в соответствии с приведённой
  схемой алгоритма.
\end{itemize}

Описание каждого алгоритма должно включать:

\begin{itemize}
\tightlist
\item
  функциональное назначение алгоритма;
\item
  входные и выходные данные (результаты выполнения);
\item
  список формальных параметров и их назначение;
\item
  пример вызова модуля или подпрограммы;
\item
  используемые технические средства;
\item
  ссылку на таблицу переменных алгоритма;
\item
  ссылку на рисунок со схемой алгоритма;
\item
  описание процесса обработки данных в соответствии со схемой;
\item
  если имеется приложение с полным текстом программы, то ссылку на
  соответствующую страницу приложения.
\end{itemize}

При описании процесса обработки данных в соответствии со схемой
алгоритма необходимо пояснить все циклы, каждую альтернативу ветвления,
принятое решение по результатам анализа альтернатив и последующие
действия.

Тексты описания алгоритмов должны быть структурными, предложения
короткими. Описание алгоритма должно отражать суть процесса обработки.

Для программы, при разработке которой использовалась
объектно-ориентированная технология, обязательно должна описываться
иерархия или диаграмма классов.

Для каждого класса желательно указать дополнительные поля и методы,
соответственно обосновывая их назначение и функции. При необходимости
здесь же можно привести алгоритмы некоторых методов.

Для пояснения особенностей реализации классов или специфики событийной
обработки можно использовать дополнительные иллюстрации, например,
диаграммы последовательности действий.

В завершении раздела описывается декомпозиция разрабатываемой программы
на модули и приводится диаграмма компоновки программного продукта.

Кроме того, в этом же разделе желательно указать вариант разработки
(«восходящая» или «нисходящая») и обосновать свой выбор.

\section{Подготовка тестовых
данных}\label{ux43fux43eux434ux433ux43eux442ux43eux432ux43aux430-ux442ux435ux441ux442ux43eux432ux44bux445-ux434ux430ux43dux43dux44bux445}

В этом разделе выбирается стратегия и методы тестирования. В
соответствии с выбранной стратегией и методами строятся примеры тестов
(обязательно с предполагаемыми результатами тестирования). Данные тестов
рекомендуется представить в виде таблиц.

\bookmarksetup{startatroot}

\chapter{Обзор инструментов моделирования. Их
установка}\label{ux43eux431ux437ux43eux440-ux438ux43dux441ux442ux440ux443ux43cux435ux43dux442ux43eux432-ux43cux43eux434ux435ux43bux438ux440ux43eux432ux430ux43dux438ux44f.-ux438ux445-ux443ux441ux442ux430ux43dux43eux432ux43aux430}

\section{Kathara}\label{kathara}

Kathara -- это бесплатное ПО с открытым исходным кодом, поддерживаемое
на всех основных операционных системах (Linux, macOs, Windows).

В среде Kathara каждое сетевое устройство реализовано в виде контейнера,
а каждое соединение эмулируется с помощью виртуальной сети.

Каждое устройство можно настроить для работы с произвольным количеством
(виртуальных) сетевых интерфейсов.

По умолчанию устройства используют образ Docker, который включает
сетевое программное обеспечение, такое как демоны маршрутизации (RIP,
OSPF и т.д.), HTTP-сервер, инструменты для работы с фаерволом
(iptables(8)) и диагностические утилиты (ping(1), traceroute(1),
tcpdump(1) и т.д.). Путем настройки соответствующего программного
обеспечения можно точно эмулировать конкретное сетевое устройство
(например, маршрутизатор).

Kathara предоставляет два альтернативных интерфейса для запуска и
настройки устройств. Набор команд с префиксом v- (vstart, vclean,
vconfig), которые позволяют запускать и управлять отдельными
устройствами, обеспечивая детальный контроль над их конфигурацией; и
набор команд с префиксом l- (lstart, lclean, linfo, lrestart, lconfig),
которые упрощают настройку предварительно сконфигурированных сетевых
сценариев, состоящих из нескольких устройств.

\subsection{Установка на
macOS}\label{ux443ux441ux442ux430ux43dux43eux432ux43aux430-ux43dux430-macos}

Для работы Kathara необходимо, чтобы был установлен и запущен Docker,
поэтому для начала обратимся к инструкции по его установке. Необходимо:
1. Скачать установщик в формате \texttt{.dmg} по следующей ссылке
https://docs.docker.com/desktop/release-notes/ 2. Дважды щелкнуть на
файл Docker.dmg, чтобы открыть установщик, затем перетащить значок
Docker в папку «Программы». По умолчанию Docker Desktop устанавливается
в папку /Applications/Docker.app. 3. Дважды щелкнуть Docker.app в папке
«Программы», чтобы запустить Docker. Установка завершена успешно.

Есть два варианта, как установить Kathara: загрузив релизный пакет с
GitHub или используя Homebrew. В первом случае необходимо: 1. Скачать
установочный файл в выбранный вами каталог из раздела Release в
официальном репозитории Kathara. 2. Запустить мастер установки и
следовать инструкциям. 3. Не забыть запустить Docker перед
использованием Kathara.

При установке через Homebrew шагов меньше. В терминале необходимо ввести
следующие команды:

\begin{minted}{shell}
brew tap KatharaFramework/kathara
brew install --cask kathara
\end{minted}

\subsection{Установка на
Linux}\label{ux443ux441ux442ux430ux43dux43eux432ux43aux430-ux43dux430-linux}

Kathara может быть установлена на всех основных дистрибутивах Linux.
Повторим, что по умолчанию Kathara использует Docker в качестве основной
среды виртуализации, однако также может работать на Kubernetes с
использованием компонента Megalos.

Для установки Docker можно воспользоваться официальной инструкцией по
установке, доступной на сайте Docker.

Рекомендуется установить эмулятор терминала xterm, который Kathara
использует по умолчанию. Команда установки зависит от используемого
дистрибутива:

\begin{verbatim}
sudo apt install xterm      # для систем на базе Debian/Ubuntu
sudo yum install xterm      # для систем на базе Red Hat/Fedora
sudo pacman -S xterm        # для Arch Linux
\end{verbatim}

\subsubsection*{Debian-based}\label{debian-based}
\addcontentsline{toc}{subsubsection}{Debian-based}

Kathara предоставляет готовые бинарные пакеты для архитектур amd64 и
arm64.

\paragraph*{Ubuntu}\label{ubuntu}
\addcontentsline{toc}{paragraph}{Ubuntu}

Для установки необходимо придерживаться следующей инструкции: 1.
Добавить открытый ключ Kathara в систему:
\texttt{sudo\ apt-key\ adv\ -\/-keyserver\ keyserver.ubuntu.com\ -\/-recv-keys\ 21805A48E6CBBA6B991ABE76646193862B759810}.
2. Подключить репозиторий Kathara:
\texttt{sudo\ add-apt-repository\ ppa:katharaframework/kathara}. 3.
Обновить список пакетов: \texttt{sudo\ apt\ update}. 4. Установить
Kathara: \texttt{sudo\ apt\ install\ kathara}.

\paragraph*{Debian 9-12}\label{debian-9-12}
\addcontentsline{toc}{paragraph}{Debian 9-12}

Kathara поддерживает все основные версии Debian (9, 10, 11 и 12) для
архитектур amd64 и arm64. Инструкция по установке во многом одинакова,
различия касаются только способа добавления ключа репозитория и версии
(кодового имени) Ubuntu, на которой основаны пакеты Kathara.

Инструкция по установке:

\begin{enumerate}
\def\labelenumi{\arabic{enumi}.}
\tightlist
\item
  Добавить открытый ключ Kathara в систему:
\end{enumerate}

\begin{itemize}
\tightlist
\item
  \textbf{Для Debian 9 и 10}:
\end{itemize}

\begin{verbatim}
sudo apt-key adv --keyserver keyserver.ubuntu.com --recv-keys 21805A48E6CBBA6B991ABE76646193862B759810
\end{verbatim}

\begin{itemize}
\tightlist
\item
  \textbf{Для Debian 11 и 12}:
\end{itemize}

\begin{verbatim}
wget -qO - "https://keyserver.ubuntu.com/pks/lookup?op=get&search=0x21805a48e6cbba6b991abe76646193862b759810" | sudo gpg --dearmor -o /usr/share/keyrings/ppa-kathara-archive-keyring.gpg
\end{verbatim}

\begin{enumerate}
\def\labelenumi{\arabic{enumi}.}
\setcounter{enumi}{1}
\tightlist
\item
  Добавить репозиторий Kathara. Можно вручную создать файл
  /etc/apt/sources.list.d/kathara.list и добавив в него строки, в
  зависимости от версии Debian \ref{tbl-ver}:
\end{enumerate}

\begin{longtable}[]{@{}
  >{\raggedright\arraybackslash}p{(\linewidth - 4\tabcolsep) * \real{0.0514}}
  >{\raggedright\arraybackslash}p{(\linewidth - 4\tabcolsep) * \real{0.0665}}
  >{\raggedright\arraybackslash}p{(\linewidth - 4\tabcolsep) * \real{0.8822}}@{}}
\caption{Содержимое файла в соответствии с версией
Debian}\label{tbl-ver}\tabularnewline
\toprule\noalign{}
\begin{minipage}[b]{\linewidth}\raggedright
\textbf{Версия Debian}
\end{minipage} & \begin{minipage}[b]{\linewidth}\raggedright
\textbf{Кодовое имя Ubuntu}
\end{minipage} & \begin{minipage}[b]{\linewidth}\raggedright
\textbf{Содержимое файла}
\end{minipage} \\
\midrule\noalign{}
\endfirsthead
\toprule\noalign{}
\begin{minipage}[b]{\linewidth}\raggedright
\textbf{Версия Debian}
\end{minipage} & \begin{minipage}[b]{\linewidth}\raggedright
\textbf{Кодовое имя Ubuntu}
\end{minipage} & \begin{minipage}[b]{\linewidth}\raggedright
\textbf{Содержимое файла}
\end{minipage} \\
\midrule\noalign{}
\endhead
\bottomrule\noalign{}
\endlastfoot
\textbf{9 / 10} & bionic &
\texttt{\textless{}br\textgreater{}deb\ http://ppa.launchpad.net/katharaframework/kathara/ubuntu\ bionic\ maindeb-src\ http://ppa.launchpad.net/katharaframework/kathara/ubuntu\ bionic\ main\textless{}br\textgreater{}} \\
\textbf{11} & focal &
\texttt{\textless{}br\textgreater{}deb\ {[}\ signed-by=/usr/share/keyrings/ppa-kathara-archive-keyring.gpg\ {]}\ http://ppa.launchpad.net/katharaframework/kathara/ubuntu\ focal\ maindeb-src\ {[}\ signed-by=/usr/share/keyrings/ppa-kathara-archive-keyring.gpg\ {]}\ http://ppa.launchpad.net/katharaframework/kathara/ubuntu\ focal\ main\textless{}br\textgreater{}} \\
\textbf{12} & jammy &
\texttt{\textless{}br\textgreater{}deb\ {[}\ signed-by=/usr/share/keyrings/ppa-kathara-archive-keyring.gpg\ {]}\ http://ppa.launchpad.net/katharaframework/kathara/ubuntu\ jammy\ maindeb-src\ {[}\ signed-by=/usr/share/keyrings/ppa-katharaframework/kathara/ubuntu\ jammy\ main\textless{}br\textgreater{}} \\
\end{longtable}

Альтернативно, можно выполнить соответствующие команды: - \textbf{Debian
9 / 10}

\begin{verbatim}
echo "deb http://ppa.launchpad.net/katharaframework/kathara/ubuntu bionic main" | sudo tee /etc/apt/sources.list.d/kathara.list
echo "deb-src http://ppa.launchpad.net/katharaframework/kathara/ubuntu bionic main" | sudo tee -a /etc/apt/sources.list.d/kathara.list
\end{verbatim}

\begin{itemize}
\tightlist
\item
  \textbf{Debian 11 / 12} (Для Debian 12 заменить focal на jammy)
\end{itemize}

\begin{verbatim}
echo "deb [ signed-by=/usr/share/keyrings/ppa-kathara-archive-keyring.gpg ] http://ppa.launchpad.net/katharaframework/kathara/ubuntu focal main" | sudo tee /etc/apt/sources.list.d/kathara.list
echo "deb-src [ signed-by=/usr/share/keyrings/ppa-kathara-archive-keyring.gpg ] http://ppa.launchpad.net/katharaframework/kathara/ubuntu focal main" | sudo tee -a /etc/apt/sources.list.d/kathara.list
\end{verbatim}

\begin{enumerate}
\def\labelenumi{\arabic{enumi}.}
\setcounter{enumi}{2}
\tightlist
\item
  Обновить список пакетов: \texttt{sudo\ apt\ update}.
\item
  Установить Kathara: \texttt{sudo\ apt\ install\ kathara}.
\end{enumerate}

\subsubsection*{Kali Linux}\label{kali-linux}
\addcontentsline{toc}{subsubsection}{Kali Linux}

Для установки на Kali Linux необходимо придерживаться инструкций для
Debian 12.

\subsubsection*{Red Hat-based}\label{red-hat-based}
\addcontentsline{toc}{subsubsection}{Red Hat-based}

\begin{enumerate}
\def\labelenumi{\arabic{enumi}.}
\tightlist
\item
  Загрузите пакет Kathara в формате \texttt{.rpm} из раздела Release
  официального репозитория на GitHub.
\item
  Перейдите в каталог, где находится загруженный файл, и выполните
  команду:
  \texttt{sudo\ rpm\ -U\ \textless{}KATHARA\_RPM\_FILE\textgreater{}.rpm}.
\item
  Перед запуском Kathara необходимо отключить \texttt{firewalld}, чтобы
  обеспечить корректную работу эмуляции сети:
  \texttt{sudo\ systemctl\ stop\ firewalld}.
\end{enumerate}

\subsubsection*{Arch-based}\label{arch-based}
\addcontentsline{toc}{subsubsection}{Arch-based}

Установка возможна двумя способами: из готового бинарного пакета или из
исходников через AUR (Arch User Repository). 1. Установка из бинарного
файла 1. Скачайте файл Kathara в формате \texttt{.pkg.tar.zst} из
раздела Release. 2. Перейдите в каталог с файлом и выполните команду:
\texttt{sudo\ pacman\ -U\ \textless{}KATHARA\_PKG\_FILE\textgreater{}.pkg.tar.zst}.
2. Установка из AUR (исходные файлы) 1. Клонируйте репозиторий Kathara
из AUR: \texttt{git\ clone\ https://aur.archlinux.org/kathara.git}. 2.
Перейдите в созданный каталог: \texttt{cd\ kathara}. 3. Соберите и
установите пакет: \texttt{makepkg\ -si}. Эта команда автоматически
выполнит сборку пакета из исходников и установит его в систему.

\subsection{Установка на
Windows}\label{ux443ux441ux442ux430ux43dux43eux432ux43aux430-ux43dux430-windows}

Kathara зависит от Docker. Если он у вас уже установлен, то можно
переходить сразу к установке эмулятора.

Скачать Docker Desktop можно с официального сайта
https://docs.docker.com/desktop/release-notes/ (Docker Desktop
Installer.exe). При установке требуется указать
\texttt{Use\ WSL\ 2\ instead\ of\ Hyper-V}, так как Kathara по умолчанию
работает с использованием Linux контейнеров. После перезапуска
компьютера проверить корректность установки можно из powershell:

\begin{minted}{powershell}
docker --version
docker run hello-world
\end{minted}

Необходимо установить Docker Desktop версии \textgreater= 2.5.0.0.

Перед использованием Kathará требуется запустить Docker. Для установки
эмулятора необходимо: 1. Скачать установочный файл в выбранную вами
директорию со страницы Releases
https://github.com/KatharaFramework/Kathara/releases. 2. Запустить
мастер установки и следовать инструкциям.

Kathara не используется из классической командной строки.

\subsection{Megalos}\label{megalos}

Megalos представляет собой расширение фреймворка Kathara, которое
добавляет поддержку Kubernetes в качестве бэкенд-системы виртуализации.
Данная функциональность была официально представлена в версии Kathara
3.0.0.

Megalos отличается от стандартного Kathara. В нем ограничено именование:
длина имен collision domain ограничена 4 символами. Сетевые интерфейсы
именуются как netX вместо ethX. Каждое устройство имеет служебный
интерфейс eth0 для коммуникации с Kubernetes, поэтому этот интерфейс
запрещено изменять. Интерфейс eth0 должен быть отключен в демонах
маршрутизации. Требуется блокировка подсети Kubernetes Pods в таблицах
маршрутизации.

Установка Megalos CNI в кластере Kubernetes требует: 1. Настроить
подключения к Kubernetes API Server через kathara settings. 2. Указать
URL и API Key для удаленного подключения.

\section{Containerlab}\label{containerlab}

Containerlab -- это инструмент, позволяющий управлять сетевыми
лабораториями на основе контейнеров. Принцип работы следующий: он
запускает контейнеры, создает виртуальное соединение между ними для
создания топологий лабораторий и управляет жизненным циклом эмуляции.

Инструмент подходит для создания произвольных linux-контейнеров, в
которых могут размещаться сетевые приложения и виртуальные функции,
также Containerlab может выступать тестовым клиентом. При всем этом он
предоставляет единый интерфейс IaaC (Infrastructure as a Code) для
управления лабораториями, который может охватывать все необходимые
варианты узлов.

\subsection{Установка на
Linux}\label{ux443ux441ux442ux430ux43dux43eux432ux43aux430-ux43dux430-linux-1}

Containerlab распространяется как Linux-пакет deb/rpm/apk для архитектур
amd64 и arm64 и может быть установлен на любой дистрибутив Debian- или
RHEL-подобного типа за несколько секунд. Для успешной работы
Containerlab необходимо заранее: - Установить Docker или Linux сервер; -
Иметь права sudo для запуска Containerlab; - Загрузить образы
контейнеров (Containerlab попытается загрузить образы во время
выполнения, если они не существуют локально). Самый простой способ
начать работу с Containerlab - использовать скрипт быстрой установки
quick-setup.sh. Он написан на bash и автоматически определяет тип ОС
(Debian, Ubuntu, RHEL, CentOS, Fedora, Rocky), после чего выполняет
установку нужных пакетов и зависимостей.

В начале скрипт проверяет: - что пользователь имеет права sudo
(требуются для установки системных пакетов); - что выполняется на
поддерживаемой Linux-системе; - наличие интернет-соединения для загрузки
зависимостей.

Далее при необходимости устанавливается Docker, Containerlab и
инструмент GitHub Cli (gh).

Чтобы установить все компоненты сразу, выполните следующую команду на
любой из поддерживаемых ОС:

\begin{verbatim}
curl -sL https://containerlab.dev/setup | sudo -E bash -s "all"
\end{verbatim}

По умолчанию это также настроит sshd в системе, чтобы неизвестные ключи
не блокировали попытки подключения по SSH. Это поведение можно
отключить, установив переменную окружения SETUP\_SSHD в значение
\texttt{false} перед запуском приведённой выше команды. Переменная
окружения устанавливается и экспортируется следующей командой:

\begin{verbatim}
export SETUP_SSHD="false"
\end{verbatim}

Чтобы завершить установку и включить выполнение команд \texttt{docker}
без sudo, выполните \texttt{newgrp\ docker} или выйдите из системы и
войдите снова.

Containerlab также настраивается для работы без sudo, и пользователь,
выполняющий скрипт быстрой установки, автоматически получает доступ к
привилегированным командам.

Чтобы установить отдельный компонент, укажите имя функции в качестве
аргумента скрипта. Например, чтобы установить только docker:

\begin{verbatim}
curl -sL https://containerlab.dev/setup | sudo -E bash -s "install-docker"
\end{verbatim}

Выйдите из текущей сессии пользователя и войдите снова, чтобы увидеть
новое двухстрочное приглашение в действии:

\begin{verbatim}
[*]─[clab]─[~]
└──>
\end{verbatim}

Если вам не надо устанавливать дополнительные инструменты, а необходим
лишь один Containerlab можно воспользоваться скриптом установки.
Например, чтобы скачать и установить последнюю версию:

\begin{verbatim}
bash -c "$(curl -sL https://get.containerlab.dev)"
\end{verbatim}

Также можно установить официальные выпуски containerlab через публичный
репозиторий APT/YUM.

\textbf{APT}

\begin{verbatim}
echo "deb [trusted=yes] https://netdevops.fury.site/apt/ /" | \
sudo tee -a /etc/apt/sources.list.d/netdevops.list

sudo apt update && sudo apt install containerlab
\end{verbatim}

\textbf{YUM}

\begin{verbatim}
sudo yum-config-manager --add-repo=https://netdevops.fury.site/yum/ && \ [](https://containerlab.dev/install/#__codelineno-9-2)echo "gpgcheck=0" | sudo tee -a /etc/yum.repos.d/netdevops.fury.site_yum_.repo [](https://containerlab.dev/install/#__codelineno-9-3)[](https://containerlab.dev/install/#__codelineno-9-4)sudo yum install containerlab
\end{verbatim}

Пакетный установщик поместит бинарный файл containerlab в каталог
\texttt{/usr/bin}, а также создаст символическую ссылку
\texttt{/usr/bin/clab\ -\textgreater{}\ /usr/bin/containerlab}. Эта
ссылка позволяет пользователям экономить на наборе команд при
использовании containerlab: clab . Containerlab также настраивается для
работы без sudo, и текущий пользователь (даже если менеджер пакетов был
вызван через sudo) автоматически получает доступ к привилегированным
командам Containerlab.

\subsection{Установка с помощью
контейнера}\label{ux443ux441ux442ux430ux43dux43eux432ux43aux430-ux441-ux43fux43eux43cux43eux449ux44cux44e-ux43aux43eux43dux442ux435ux439ux43dux435ux440ux430}

Containerlab также доступен в виде образа контейнера. Последнюю версию
можно загрузить командой:

\begin{verbatim}
docker pull ghcr.io/srl-labs/clab
\end{verbatim}

Чтобы выбрать любую из выпущенных версий начиная с \textbf{0.19.0},
используйте номер версии в качестве тега, например:

\begin{verbatim}
docker pull ghcr.io/srl-labs/clab:0.19.0
\end{verbatim}

Загруженный контрейнер запускается командой:

\begin{verbatim}
docker run --rm -it --privileged \
    --network host \
    -v /var/run/docker.sock:/var/run/docker.sock \
    -v /var/run/netns:/var/run/netns \
    -v /etc/hosts:/etc/hosts \
    -v /var/lib/docker/containers:/var/lib/docker/containers \
    --pid="host" \
    -v $(pwd):$(pwd) \
    -w $(pwd) \
    ghcr.io/srl-labs/clab bash
\end{verbatim}

Рассмотрим основные требуемые привилегии:

\begin{itemize}
\tightlist
\item
  \texttt{-\/-privileged}: полный доступ к хосту для создания и
  управления сетевым пространством;
\item
  \texttt{-\/-network\ host}: доступ к сетевому стеку хоста;
\item
  \texttt{-\/-pid="host"}: доступ к пространству процессов хоста.
\end{itemize}

В запущенном контейнере можно использовать те же команды
\textbf{deploy/destroy/inspect} для управления лабораториями. Можно
развернуть лабораторию альтернативным способом без запуска оболочки,
например:

\begin{verbatim}
docker run --rm -it --privileged \
# <параметры опущены>
-w $(pwd) \
ghcr.io/srl-labs/clab deploy -t somelab.clab.yml
\end{verbatim}

\subsection{Ручная
установка}\label{ux440ux443ux447ux43dux430ux44f-ux443ux441ux442ux430ux43dux43eux432ux43aux430}

Если дистрибутив Linux не поддерживает установку deb/rpm пакетов,
containerlab можно установить из архива:

\begin{verbatim}
LATEST=$(curl -s https://github.com/srl-labs/containerlab/releases/latest | \
       sed -e 's/.*tag\/v\(.*\)\".*/\1/')
curl -L -o /tmp/clab.tar.gz "https://github.com/srl-labs/containerlab/releases/download/v${LATEST}/containerlab_${LATEST}_Linux_amd64.tar.gz"
mkdir -p /etc/containerlab
tar -zxvf /tmp/clab.tar.gz -C /etc/containerlab
mv /etc/containerlab/containerlab /usr/bin && chmod a+x /usr/bin/containerlab
\end{verbatim}

\subsection{Обновление}\label{ux43eux431ux43dux43eux432ux43bux435ux43dux438ux435}

Чтобы обновить containerlab до последней доступной версии, выполните:

\begin{verbatim}
sudo -E containerlab version upgrade
\end{verbatim}

Эта команда скачает скрипт установки и обновит инструмент до последней
версии.

\subsection{Установка с помощью
исходников}\label{ux443ux441ux442ux430ux43dux43eux432ux43aux430-ux441-ux43fux43eux43cux43eux449ux44cux44e-ux438ux441ux445ux43eux434ux43dux438ux43aux43eux432}

Чтобы собрать containerlab из исходников есть два варианта:

\begin{itemize}
\tightlist
\item
  \textbf{go build};
\item
  \textbf{goreleaser}. В первом случае, требуется склонировать
  репозиторий, перейти в корневой каталог и выполнить команду
  \texttt{go\ build}:
\end{itemize}

\begin{verbatim}
git clone https://github.com/srl-labs/containerlab
cd containerlab
go build
\end{verbatim}

Во втором случае, для сборки запускается:

\begin{verbatim}
goreleaser --snapshot --skip-publish --rm-dist
\end{verbatim}

\subsection{Удаление}\label{ux443ux434ux430ux43bux435ux43dux438ux435}

Отдельно рассмотрим удаление инструмента. Чтобы удалить containerlab,
установленный через скрипт или пакеты:

\begin{itemize}
\tightlist
\item
  \textbf{Для Debian-систем:}
\end{itemize}

\begin{verbatim}
apt remove containerlab
\end{verbatim}

\begin{itemize}
\tightlist
\item
  \textbf{Для RPM-систем:}
\end{itemize}

\begin{verbatim}
yum remove containerlab
# or
dnf remove containerlab
\end{verbatim}

\subsection{Установка на
Windows}\label{ux443ux441ux442ux430ux43dux43eux432ux43aux430-ux43dux430-windows-1}

Для того, чтобы развернуть Containerlab на Windows, требуется
использовать WSL
(~\href{https://learn.microsoft.com/en-us/windows/wsl/}{Windows
Subsystem Linux}). WSL позволяет пользователям запускать легковесные
Linux VM внутри Windows, это можно использовать для развертывания
Containerlab.

Если WSL не установлен, исправить это можно одной командой:

\begin{minted}{powershell}
wsl --install
\end{minted}

Далее расcмотрим установку Containerlab. В Powersher от имени
администратора включаем компоненты Windows, необходимые для работы
подсистемы Linux: активируем базовую функциональность WSL и платформу
виртуальных машин для использования более производительной версии WSL 2:

\begin{minted}{powershell}
dism.exe /online /enable-feature /featurename:Microsoft-Windows-Subsystem-Linux /all /norestart

dism.exe /online /enable-feature /featurename:VirtualMachinePlatform /all /norestart
\end{minted}

После выполнения этих команд и~перезагрузки~компьютера, нужно с помощью
команды~\texttt{wsl\ -\/-set-default-version\ 2}~задать WSL 2 как версию
по умолчанию.

Загрузим файл \texttt{.wsl} со
\href{https://github.com/srl-labs/wsl-containerlab/releases/tag/0.69.3-1.0}{страницы
релизов}. Просто дважды щелкнем по файлу, и дистрибутив будет
установлен. Альтернативно его можно установить с
помощью\texttt{wsl\ -install\ -from-file\ C:\textbackslash{}path\textbackslash{}to\textbackslash{}clab.wsl}.
Containerlab должен появиться как программа в меню \enquote{Пуск}.
Запустим Containerlab:

\begin{minted}{powershell}
wsl -d Containerlab
\end{minted}

При первом запуске Containerlab WSL выполняется начальный этап
настройки, который помогает настроить дистрибутив. Настраивается
оболочка и командная строка. После этого SSH-ключи будут скопированы или
сгенерированы (если ключей не существует). Это необходимо для
обеспечения доступа к SSH без пароля, что улучшает интеграцию с DevPod.

\subsection{Установка на
MacOS}\label{ux443ux441ux442ux430ux43dux43eux432ux43aux430-ux43dux430-macos-1}

Запуск containerlab на macOS возможен как на ARM (M1/M2/M3 и т.д.), так
и на Intel-чипах. Долгое время существовали ограничения для ARM-чипов,
но с появлением ARM64-совместимых сетевых ОС, таких как Nokia SR Linux и
Arista cEOS, а также благодаря эмуляции Rosetta для x86\_64-систем,
теперь можно запускать containerlab и на Mac с ARM чипом.

Для работы Containerlab на macOS используется Linux-виртуальная среда,
поскольку macOS не имеет нативной поддержки Linux-сетевых пространств. В
качестве такой среды рекомендуется использовать OrbStack --- лёгкий
гипервизор, совместимый с Docker и Linux-ядром.

\subsubsection*{Порядок
установки}\label{ux43fux43eux440ux44fux434ux43eux43a-ux443ux441ux442ux430ux43dux43eux432ux43aux438}
\addcontentsline{toc}{subsubsection}{Порядок установки}

\begin{enumerate}
\def\labelenumi{\arabic{enumi}.}
\tightlist
\item
  \textbf{Установить OrbStack}. Скачать и установить OrbStack с
  официального сайта \url{https://orbstack.dev}. OrbStack создаёт лёгкую
  виртуальную машину Linux, в которой работает Docker и все сетевые
  инструменты, необходимые Containerlab.
\item
  \textbf{Запустить Linux-среду внутри OrbStack}. После установки
  открыть терминал OrbStack и войти в Linux-оболочку (по умолчанию
  Ubuntu-подобная система).
\item
  \textbf{Установить Docker и Containerlab}. Внутри OrbStack-Linux
  выполнить стандартную установку для Linux:
  \texttt{curl\ -sL\ https://containerlab.dev/setup\ \textbar{}\ sudo\ -E\ bash\ -s\ "all"}.
  Эта команда установит Docker, Docker Compose, GitHub CLI и сам
  Containerlab.
\end{enumerate}

!{[}{[}Screenshot 2025-10-22 at 18.12.49.png{]}{]}

\begin{enumerate}
\def\labelenumi{\arabic{enumi}.}
\setcounter{enumi}{3}
\tightlist
\item
  \textbf{Проверить установку}. После завершения выйти и войти снова,
  чтобы активировать группы пользователей, и проверить:
\end{enumerate}

\begin{verbatim}
clab version
docker ps
\end{verbatim}

!{[}{[}Screenshot 2025-10-22 at 18.15.15.png{]}{]}

\begin{enumerate}
\def\labelenumi{\arabic{enumi}.}
\setcounter{enumi}{4}
\tightlist
\item
  \textbf{Проверить архитектуру сетевых образов}. На ARM-Mac
  рекомендуется использовать контейнерные образы, собранные под
  архитектуру arm64. Проверить это можно командой:
\end{enumerate}

\begin{verbatim}
docker image inspect <image> -f '{{.Architecture}}'
\end{verbatim}

\begin{enumerate}
\def\labelenumi{\arabic{enumi}.}
\setcounter{enumi}{5}
\tightlist
\item
  \textbf{Развернуть лабораторию}. Теперь можно развернуть тестовую
  лабораторию, например:
  \texttt{sudo\ clab\ deploy\ -t\ examples/srl01.clab.yml}.
\end{enumerate}

!{[}{[}Screenshot 2025-10-22 at 18.25.40.png{]}{]}

Зайдем на первый хост, посмотрим существующие интерфейсы и поднимем
интерфейс ethernet-1/1, добавим ipv4 адрес 10.0.0.1/31.

!{[}{[}Screenshot 2025-10-22 at 18.34.40.png{]}{]}

Проверка добавления интерфейса.

!{[}{[}Screenshot 2025-10-22 at 18.34.53.png{]}{]}

Аналогично делаем на втором узле. Присваиваем ipv4 адрес 10.0.0.2/31.

!{[}{[}Screenshot 2025-10-22 at 18.37.52.png{]}{]}

!{[}{[}Screenshot 2025-10-22 at 18.38.14.png{]}{]}

Пропингуем первый хост со второго: ping 10.0.0.1

С помощью команды \texttt{containerlab\ graph\ -t\ srl02.clab.yml}
посмотрим на графическое изображение сети. !{[}{[}Screenshot 2025-10-22
at 19.08.04.png{]}{]}

\bookmarksetup{startatroot}

\chapter*{Заключение}\label{ux437ux430ux43aux43bux44eux447ux435ux43dux438ux435}
\addcontentsline{toc}{chapter}{Заключение}

\markboth{Заключение}{Заключение}

Заключение курсовой работы является ее завершающей частью, которая
содержит выводы и предложения, полученные в ходе исследования, авторскую
оценку работы с точки зрения решения задач, поставленных в курсовой
работе. Объем заключения --- 1--2 страницы.

Список использованных источников должен содержать сведения о нормативных
правовых актах, учебных, методических и научных изданиях (на русском и
иностранном языках), публикациях в периодической печати, а также базах
данных, информационно-справочных системах и Интернет-ресурсах,
использованных студентом в ходе выполнения курсовой работы. Для
написания курсовой работы студент должен использовать не менее 30
источников.

\bookmarksetup{startatroot}

\chapter*{Список используемых
сокращений}\label{ux441ux43fux438ux441ux43eux43a-ux438ux441ux43fux43eux43bux44cux437ux443ux435ux43cux44bux445-ux441ux43eux43aux440ux430ux449ux435ux43dux438ux439}
\addcontentsline{toc}{chapter}{Список используемых сокращений}

\markboth{Список используемых сокращений}{Список используемых
сокращений}

\section*{Англоязычные
сокращения}\label{ux430ux43dux433ux43bux43eux44fux437ux44bux447ux43dux44bux435-ux441ux43eux43aux440ux430ux449ux435ux43dux438ux44f}
\addcontentsline{toc}{section}{Англоязычные сокращения}

\markright{Англоязычные сокращения}

AQM --- Active Queue Management --- активное управление очередью

\section*{Русскоязычные
сокращения}\label{ux440ux443ux441ux441ux43aux43eux44fux437ux44bux447ux43dux44bux435-ux441ux43eux43aux440ux430ux449ux435ux43dux438ux44f}
\addcontentsline{toc}{section}{Русскоязычные сокращения}

\markright{Русскоязычные сокращения}

ПМП --- Планетарный механизм поворота

\cleardoublepage
\phantomsection
\addcontentsline{toc}{part}{Приложения}
\appendix

\chapter{Приложение}\label{ux43fux440ux438ux43bux43eux436ux435ux43dux438ux435}

Приложения помещают после списка использованных источников. В составе
приложений могут быть представлены таблицы, диаграммы, схемы, графики,
позволяющие раскрыть тему, имеющие вспомогательное значение для текста
работы, но не включенные в ее основную часть.

Название конкретного приложения должно раскрывать его содержание.

Приводятся табличные и графические результаты (рис.
\ref{fig-luneburg_rays_2}).

\begin{figure}

\centering{

\includegraphics[width=0.7\linewidth,height=\textheight,keepaspectratio]{text/../image/luneburg_rays.png}

}

\caption{\label{fig-luneburg_rays_2}Прохождение лучей через линзу
Люнеберга}

\end{figure}%

\chapter{Код
программы}\label{ux43aux43eux434-ux43fux440ux43eux433ux440ux430ux43cux43cux44b}

\begin{minted}{julia}
using Pkg
Pkg.add("DifferentialEquations")
using DifferentialEquations

pr = 0.0 # initial value of the drop function
T = 0.5 # round-trip time, RTT
N = 60.0 # number of TCP sessions
c = 10.0 # service intensity, Mbps
packet_size = 500.0 # package size, bit
C = 125000.0*c / packet_size # service intensity, packets
wq = 0.0004 # EWMS parameter
q_min = 0.25 # minimum queue threshold
q_max = 0.50 # maximum queue threshold
R = 300.0 # queue size
p_max = 0.1 # maximum drop probability
w_max = 32.0 #maximum window size
p = (T, N, C, wq, q_min, q_max, R, p_max, w_max) # vector of system parameters
\end{minted}

\chapter{Код программного комплекса расчёта
линз}\label{ux43aux43eux434-ux43fux440ux43eux433ux440ux430ux43cux43cux43dux43eux433ux43e-ux43aux43eux43cux43fux43bux435ux43aux441ux430-ux440ux430ux441ux447ux451ux442ux430-ux43bux438ux43dux437}

\section{Файл для сборки программного
комплекса}\label{ux444ux430ux439ux43b-ux434ux43bux44f-ux441ux431ux43eux440ux43aux438-ux43fux440ux43eux433ux440ux430ux43cux43cux43dux43eux433ux43e-ux43aux43eux43cux43fux43bux435ux43aux441ux430}

\begin{minted}{make}

\end{minted}

\section{Расчёт линз в полярных
координатах}\label{ux440ux430ux441ux447ux451ux442-ux43bux438ux43dux437-ux432-ux43fux43eux43bux44fux440ux43dux44bux445-ux43aux43eux43eux440ux434ux438ux43dux430ux442ux430ux445}

\begin{minted}{julia}

\end{minted}


\printbibliography



\end{document}
